\documentclass[10pt,a4paper]{amsart}
\usepackage[margin=19mm]{geometry}
\usepackage{amsmath,amssymb,mathtools}
\usepackage[scr=rsfs,cal=euler]{mathalfa}
\usepackage{xfrac}
\usepackage[unicode,hidelinks,bookmarks=false]{hyperref}
\newcommand{\ie}{i.e.\ }
\newcommand{\cf}{cf.\ }
\newcommand{\resp}{respectively\ }
\newcommand{\Subsection}{\subsection}
\providecommand{\nobreakdash}{\nobreak\hbox{-}}
\setlength{\emergencystretch}{2em}
\multlinegap0pt
\usepackage[shortlabels]{enumitem}
\usepackage{amssymb}
\usepackage{mathtools}
\usepackage{xcolor}
\newtheorem{proposition}{Proposition}
\title{On Minimal Noncommutative Rings}
\author{V. V. Bavula, N. Blacher}
\date{Source: arXiv:2608.09370v1 (CC BY 4.0)}
\begin{document}
\maketitle

\noindent\emph{Source and licence.} On Minimal Noncommutative Rings. Version arXiv:2608.09370v1 (CC BY 4.0), distributed by arXiv under the Creative Commons Attribution 4.0 International licence, http://creativecommons.org/licenses/by/4.0/, as recorded in the arXiv licence metadata for this exact version and shown on the abstract page https://arxiv.org/abs/2608.09370v1. Full licence notice: \texttt{licenses/arxiv-2608.09370-CC-BY-4.0.txt}.\par\smallskip

\noindent\textit{Editorial scope.} Counted unit: the article's proposition FirstTwoClasses (source0.tex lines 128--134). The article's complete original proof of it is delivered with it (source0.tex lines 136--194, one \texttt{proof} environment of 59 lines). No mathematics is written, repaired or re-derived: the statement and the proof are the article's own lines, in the article's own notation, and no retained line is substituted or re-flowed.  No further source-local context is needed: every cross-reference of every retained line resolves inside the two delivered spans.  Nothing was omitted from any retained span, so the package carries no omission record.  No retained line contains a citation, so the wrapper carries no bibliography and no bibliography entry.  No retained line contains a figure environment, an image-inclusion command or drawing code, so no image is delivered and the package declares no asset dependency.  Editorial formatting only, in addition: an AMS \texttt{amsart} preamble that restates the article's own declarations and macros used by the retained lines, namely the environments \texttt{proposition} on separate counters, together with the standard packages those lines need.  The retained environments print in this document as Proposition 1 in order of appearance, whereas the article prints the same items with the numbering of its own full document.  The complete native source of the article as distributed in the arXiv e-print for this exact version is bundled unchanged as \texttt{sources/207237/source0.tex}.
\begin{proposition}\label{FirstTwoClasses}
For any prime $p$, the following hold.
\begin{itemize}[topsep=0pt]
\item[(1)] The ring $U_p$ is minimal noncommutative.
\item[(2)] The ring $B_{p,n,i}$ is minimal noncommutative if and only if $n=q^m$ and $i=jq^{m-1}$ for some prime $q$ and integers $m\geq 1$ and $1\leq j\leq q-1$.
\end{itemize}
\end{proposition}

\begin{proof}
Checking (1) is an easy exercise. We turn to (2) and first prove the condition is necessary. We will show that $B_{p,n,i}$ has a noncommutative proper subring if either $n$ is not a prime power, or $n=q^m$ but $i$ is not of the form $jq^{m-1}$. In either case, there is a natural number $1\leq k<n$ that divides $n$ but does not divide $i$. In particular, the field $F_{p^n}$ contains a proper subfield $F_{p^k}$. Then the ring $B_{p,k,i}=\left\lbrace\begin{pmatrix}
x^{p^i}&y\\
0&x
\end{pmatrix}
:x,y\in F_{p^k}\right\rbrace$ is a proper subring of $B_{p,n,i}$. We will show that $B_{p,k,i}$ is not commutative. Considering the commutator
\[\left[\begin{pmatrix}
x^{p^i}&0\\
0&x
\end{pmatrix},\begin{pmatrix}
0&1\\
0&0
\end{pmatrix}\right]=\begin{pmatrix}
0&x^{p^i}-x\\
0&0
\end{pmatrix}\]
it suffices to show that there is an element $x\in F_{p^k}$ such that $x^{p^i}\neq x$. Since $k$ does not divide $i$ we can write $i=\alpha k+\beta$ with $0<\beta<k$. Pick $x\in F_{p^k}$ such that $x^{p^\beta}\neq x$. Such an $x$ exists since the lowest degree monic polynomial identity in one variable satisfied by $F_{p^k}$ is $X^{p^k}-X$. Now
\[x^{p^{\alpha k}}=\left(x^{p^k}\right)^{p^{(\alpha-1)k}}=x^{p^{(\alpha-1)k}}=\cdots=x^{p^k}=x\]
and therefore
\[x^{p^i}=x^{p^{\alpha k+\beta}}=x^{p^{\alpha k}p^\beta}=\left(x^{p^{\alpha k}}\right)^{p^\beta}=x^{p^\beta}\neq x\]
as required.

Now we turn to the opposite implication. That is we show that $R=B_{p,n,i}$ is a minimal noncommutative ring when $n=q^m$ is a prime power and $i$ is of the form $jq^{m-1}$ for $1\leq j\leq q-1$. That $R$ is not commutative follows as above from the fact that $X^{p^i}-X$ is not an identity for the field $F_{p^n}$. We start by showing that every proper homomorphic image $R/I$ is commutative. Since every commutator in $R$ is a strictly upper triangular matrix, it is enough to show that every matrix of the form $\begin{pmatrix}
0&z\\
0&0
\end{pmatrix}$ is contained in $I$. Pick a nonzero element $\begin{pmatrix}
x^{p^i}&y\\
0&x
\end{pmatrix}$ of $I$. If $x=0$ then $y\neq 0$ and $I$ contains
\[\begin{pmatrix}
0&y\\
0&0
\end{pmatrix}\begin{pmatrix}
\left(y^{-1}z\right)^{p^i}&0\\
0&y^{-1}z
\end{pmatrix}=\begin{pmatrix}
0&z\\
0&0
\end{pmatrix}\]
as required. If $x\neq 0$ then $I$ contains
\[\begin{pmatrix}
x^{p^i}&y\\
0&x
\end{pmatrix}\begin{pmatrix}
0&x^{-p^i}z\\
0&0
\end{pmatrix}=\begin{pmatrix}
0&z\\
0&0
\end{pmatrix}\]
as required. Thus $I$ contains $[R,R]$, i.e. $R/I$ is commutative. It remains to prove that every proper subring of $R$ is commutative. Let $S$ be a subring of $R$. Consider the homomorphism $\pi:S\to F_{p^n}$ mapping each matrix in $S$ to its bottom-right entry. The image of $\pi$ must be a subfield $F_{p^k}$ of $F_{p^n}$, while its kernel must be an $F_{p^k}$-subspace of $F_{p^n}$.

First we consider the case when $\pi(S)=F_{p^n}$. Then either $\ker(\pi)=0$ and $S$ is commutative, or $\ker(\pi)\cong F_{p^n}$ and $S=R$. So we can now assume that $\pi(S)=F_{p^k}$ is a proper subfield of $F_{p^n}$. Since $n=q^m$, this means $k=q^l$ with $0\leq l\leq m-1$. In this case, since $i=jq^{m-1}$ we have $i=kh$ with $h\geq 1$. Therefore for any $x\in F_{p^k}$ we have
\[x^{p^i}=x^{p^{kh}}=\left(x^{p^k}\right)^{p^{k(h-1)}}=x^{p^{k(h-1)}}=\cdots=x^{p^k}=x.\]
Thus every element of $S$ is of the form $\begin{pmatrix}
x&y\\
0&x
\end{pmatrix}$ and it follows that $S$ is commutative, completing the proof.
\end{proof}

\end{document}
